The urgent demand for sustainable energy generation and environmental protection has accelerated interest in solar-driven photocatalysis as a practical route to mitigate these global challenges.\cite{Tong2025RecentPhotocatalysis} Semiconductor-based photocatalysts, capable of harvesting sunlight to drive chemical reactions such as water splitting for hydrogen production, carbon dioxide reduction, and degradation of organic pollutants, provide a clean and scalable approach towards a clean energy future.\cite{Tofil2025ModifiedReview} Among the wide range of semiconductor photocatalysts investigated, graphitic carbon nitride (g-C$_3$N$_4$) has emerged as a promising two-dimensional (2D) metal-free polymeric photocatalyst due to its suitable bandgap for visible light activity, chemical stability, and ease of synthesis, thereby making it highly relevant for various sustainable applications.\cite{Fang2025VersatileApplications,Khan2025RecentReview}

Graphitic carbon nitride exists in two structural frameworks, comprising either s-triazine (C$_3$N$_3$) units or tri-s-triazine (C$_6$N$_7$, also known as heptazine) units. Both structures consist of fused 6-membered rings with sp$^2$ hybridized C and N atoms, forming delocalized $\pi$-conjugated polymeric networks. 
Moreover, as revealed by x-ray diffraction (XRD, Cu target), g-C$_3$N$_4$ is a hexagonal layered structure with two characteristic peaks at approximately 2$\theta$ of 13.04$^\circ$ and 27.6$^\circ$, attributed to the $\langle 10\bar{1}0 \rangle$ in-plane periodicity and $\langle 0002 \rangle$ interlayer stacking governed by weak van der Waals interactions, respectively. While s-triazine and heptazine share a layered graphitic arrangement, heptazine units, consisting of three fused triazine rings, are energetically more favourable.\cite{Zhao2023Two-dimensionalProcesses,PhuongLy2025DecipheringWater,Rizwan2025RoleApplications} Zhang \textit{et al}.\cite{Zhang2025FunctionalNitride} report the structural evolution from triazine to heptazine during thermal polymerization and highlights new theoretical insights for synthesizing heptazine-based g-C$_3$N$_4$ at lower temperatures.

Unlike graphene, which consists mainly of carbon atoms, polymeric g-C$_3$N$_4$ contains both carbon and nitrogen atoms, resulting in increased structural complexity. Consequently, uncertainty regarding the precise structure of polymeric g-C$_3$N$_4$ and the identification of its true ground-state structure continues to present significant challenges for both experimental and theoretical investigations. For example, the conventional synthesis route of polymeric carbon nitride, such as thermal polycondensation in air, typically yields a semicrystalline structure with mixed motifs, thereby complicating precise structural characterization.\cite{Negro2023CombinedProperties}
The reduced crystallinity of g-C$_3$N$_4$ is commonly attributed to incomplete polymerization and the persistence of hydrogen bonding within the synthesized sample, which results in a partially condensed framework of a melon-type g-C$_3$N$_4$ structure.\cite{Im2023StructureSimulations}

Recently, Negro \textit{et al}.\cite{Negro2023CombinedProperties} conducted a combined experimental and theoretical investigation on g-C$_3$N$_4$ structure using different characterization techniques such as XRD, FTIR, SEM, UV-Vis, and data from density functional theory tools. Their study demonstrated that thermal polymerization of melamine produces a mixture of structural motifs comprising less-condensed melon frameworks together with highly condensed g-C$_3$N$_4$ domains. Another study by Zhang \textit{et al}.\cite{Zhang2023Salt-MediatedPhotocatalysis} provides insights into the structural transformation mechanism of g-C$_3$N$_4$ synthesized via a molten salt-assisted procedure. They report a sequential phase evolution in which melon converts first into poly(heptazine imide) as an intermediate phase, and subsequently into poly(triazine imide) as the more stable phase. Collectively, earlier works have demonstrated that the structure, and consequently the properties, of g-C$_3$N$_4$ can be modulated via different synthesis routes such as precursor composition, condensation temperature, reaction rates and related factors.\cite{Ma2024TheModification,Zhang2025FunctionalNitride} Nonetheless, the fundamental question of what is the ground state structure of fully polymerized heptazine-based g-C$_3$N$_4$ has not been answered through experimental investigation. 

In an effort to resolve the true ground state of pure heptazine-based  g-C$_3$N$_4$, numerous theoretical investigations based on first-principles density functional theory (DFT) and molecular dynamics have been conducted.\cite{Wang2018TheStudy,Wang2017DeterminationValidation,Gao2020CorrugationStudies}
While numerous studies employ a graphene-like planar model of $P\bar{6}m2$ symmetry for the g-C$_3$N$_4$ layer structure, other studies have allowed the geometry to flex into nonplanar structures.
Recent theoretical work\cite{Im2023StructureSimulations,Fao2022TheoreticalPd,Hartley2021ComputationalSplitting} has demonstrated that out-of-plane structures, such as buckled or corrugated layers, are energetically more favourable and offer a more realistic model of the fully polymerized heptazine phase. Moreover, Zhao \textit{et al}. \cite{zhao2021NewNitride} proposed three corrugated structures of $P321$, $P3m1$ and $Pca2_1$ symmetry for experimentally sythesized g-C$_3$N$_4$. These were found to be energicatlly more stable than the conventional flat $P\bar{6}m2$ model. Other theoretical investigations  that adopted a flat g-C$_3$N$_4$ model as the initial structure for adsorption or dopant studies consistently yielded buckled configurations, with planarity not being preserved once the symmetry is broken in the calculation.\cite{Azofra2016AConversion,Zhou2022TheStudy,Zhang2017PhotocatalyticG-C3N4b,Ji2016Heptazine-basedMembrane} These observations collectively suggest that out-of-plane buckling is an intrinsic feature of g-C$_3$N$_4$. However, it remains challenging, experimentally, to detect the subtly buckled heptazine phase.
When the basal plane of the layered structure is viewed from above
(with atomic-resolution microscopy techniques) the corrugated and planar g-C$_3$N$_4$ structures would be virtually indistinguishable.

Another important consideration in the search for the true ground state of heptazine-based g-C$_3$N$_4$  requires careful evaluation of how the two-dimensional (2D) layers are arranged to form a stacked three-dimensional (3D) bulk structure. The different stacking configurations between adjacent g-C$_3$N$_4$ sheets play a decisive role in determining the structural stability and electronic properties of the bulk phase. Although many theoretical studies\cite{Yang2021DFTApplications,Fu2024PotassiumInvestigation,Gorai2023InsightPrinciple} have modelled exfoliated  g-C$_3$N$_4$ as an isolated monolayer to rationalize enhanced photocatalytic activity, the persistence of the (0002) diffraction peak indicates evidence of residual interlayer stacking. As a result, interlayer interactions and $\pi$-$\pi$ overlap between adjacent sheets must not be neglected in describing the ideal phase of heptazine-based g-C$_3$N$_4$. In 2018, Sun \textit{et al.} \cite{Sun2018FirstArrangement} investigated the electronic properties of graphitic carbon nitride with different building blocks and staggered between sheets in the bulk phase. In their study, triazine- and heptazine-based structures were examined under different stacking configurations ($P\bar{6}m2$, $R3m$, $Cmcm$, and $P6_3/mmc$), and the results showed that bulk g-C$_3$N$_4$ polymorphs exhibit narrower band gaps than the corresponding monolayers, with the reduction attributed to interlayer band overlap between adjacent CN sheets. Other previous studies 
have also indicated that polymeric carbon nitride comprises stacked layers that preferentially adopt AB stacking registries as compared 
to AA.\cite{Hartley2021ComputationalSplitting,Zuluaga2015StructuralG-C3N4}
However, the precise stacking order and local structural motifs remain uncertain due to experimental constraints, motivating further investigation. 

A more recent study \cite{Im2025UnravelingTransitions} demonstrated that the electronic and optical responses of polymeric g-C$_3$N$_4$ are dependent on the microstructural features, particularly the degree of condensation, out-of-plane corrugation, and interlayer interactions. These conclusions are consistent with prior work,\cite{Im2023StructureSimulations,Negro2023CombinedProperties} providing a strong rationale to systematically investigate the effect of buckling and stacking configurations in developing a more robust theoretical model for the fully condensed heptazine phase of g-C$_3$N$_4$.

Mridha \textit{et al}. \cite{Mridha2025UsingG-C3N4} recently advanced the earlier work of Im and colleagues \cite{Im2025UnravelingTransitions} by providing a dynamic analysis of the ground-state structure of g-C$_3$N$_4$, systematically examining heptazine-, h-triazine-, and o-triazine-based polymorphs in both bulk and monolayer models. Their phonon calculations revealed that all planar polymorphs are dynamically unstable and relax to buckled structures, with heptazine identified as the energetically preferred structure. 

While earlier studies have confirmed that buckling is intrinsic to the stability of g-C$_3$N$_4$, their scope remained limited to idealized monolayers and bulk phases, without assessing the influence of stacking registries on the relative stability of corrugated heptazine structures. To address this gap, the present study systematically investigates the interplay between buckling and different stacking configurations in heptazine-based g-C$_3$N$_4$, with a direct comparison between planar/buckled monolayers and the energetically preferred corrugated bulk phase. In addition, we have investigated how altering the assumed buckling and stacking structure affects the properties of phosphorus and nickel doping of g-C$_3$N$_4$. Our results reveal that the combination of buckling and appropriate stacking configurations stabilizes heptazine by yielding lower-energy corrugated structures.


